\documentclass[11pt]{article}
\usepackage[margin=1in]{geometry}
\usepackage{amsmath,amssymb}
\usepackage{mathrsfs}
\usepackage{hyperref}

\newcommand{\om}{\omega}
\newcommand{\omone}{\omega_1}
\newcommand{\cG}{\mathscr G}
\newcommand{\cM}{\mathscr M}
\newcommand{\cX}{\mathscr X}
\newcommand{\Sz}{\mathscr{S\!Z}}

\title{Literature-Implied Answers to Two Ideal-Lattice Questions for \(C([0,\omega_1])\)}
\author{Run \texttt{fa\_banach\_001}}
\date{2026-06-23}

\begin{document}
\sloppy
\maketitle

\paragraph{Status.}
\textbf{Literature-implied answer.} The closing remark of Kania--Laustsen
\cite{KL2011} asks two questions about the lattice of closed ideals of
\(\mathscr B(C([0,\om_1]))\). Later work of Kania--Laustsen \cite{KL2014}
answers them after identifying the factor space used in the first question
and reading the improved lattice diagram in the isomorphic
\(C_0[0,\om_1)\) model.

\paragraph{Source questions.}
In the final remark of \cite{KL2011}, after constructing an Eberlein compact
space \(K\), the authors prove
\[
 \cX(C([0,\om_1]))+\overline{\cG}_{c_0(\om_1)}(C([0,\om_1]))
 \subsetneq
 \overline{\cG}_{C(K)}(C([0,\om_1]))
 \subseteq
 \cM
\]
and say that they do not know whether the final inclusion is proper, adding
that they conjecture it is. They then ask whether
\[
 \overline{\cG}_{C(K_\alpha)\oplus c_0(\om_1)}(C([0,\om_1]))
 \subseteq
 \Sz_{\alpha+1}(C([0,\om_1]))
\]
is proper for some, or each, countable ordinal \(\alpha\), where
\(K_\alpha=[0,\om^{\om^\alpha}]\).

\paragraph{Identification of the first factor space.}
The compact space \(K\) in \cite{KL2011} is the one-point compactification of
the disjoint union of \(K_\alpha=[0,\om^{\om^\alpha}]\),
\(\alpha<\om_1\). Hence
\[
 C(K)\cong
 \mathbb K\oplus
 \left(\bigoplus_{\alpha<\om_1} C(K_\alpha)\right)_{c_0}.
\]
Let \(A=\{\om^{\om^\alpha}:\alpha<\om_1\}\subset[0,\om_1)\). The non-scalar
summand is exactly \(E_A=(\oplus_{\gamma\in A} C[0,\gamma])_{c_0}\) in the
notation of \cite{KL2014}. Lemma 2.4 of \cite{KL2014} states that
\(C_0(L_0)\) is isometric to \(E_{\om_1}\), and that \(E_A\cong E_{\om_1}\)
whenever \(A\) is uncountable. The same lemma also gives enough absorption
to remove the scalar summand: \(E_{\om_1}\cong \mathbb K\oplus E_{\om_1}\).
Thus
\[
 C(K)\cong C_0(L_0).
\]

\paragraph{Answer to the conjectured proper inclusion.}
The later paper records, using the preceding work \cite{KL2011,KKL}, that the
Loy--Willis ideal on \(C_0[0,\om_1)\) is precisely the ideal of operators
factoring through \(C_0(L_0)\):
\[
  \cM=\cG_{C_0(L_0)}.
\]
It also notes that the Banach algebras
\(\mathscr B(C_0[0,\om_1))\) and \(\mathscr B(C[0,\om_1])\) have isomorphic
lattices of closed ideals because \(C_0[0,\om_1)\cong C[0,\om_1]\).
Combining these facts with \(C(K)\cong C_0(L_0)\), the source inclusion
\[
 \overline{\cG}_{C(K)}(C([0,\om_1]))\subseteq \cM
\]
is an equality. Therefore the conjecture in \cite{KL2011} that the inclusion
is proper is false.

\paragraph{Answer to the Szlenk-inclusion question.}
For every countable ordinal \(\alpha\ge 1\), \cite{KL2014} proves
\[
 \overline{\cG}_{C(K_\alpha)\oplus c_0(\om_1)}(C_0[0,\om_1))
 \subsetneqq
 \overline{\cG}_{c_0(\om_1,C(K_\alpha))}(C_0[0,\om_1))
 \subseteq
 \Sz_{\alpha+1}(C_0[0,\om_1)).
\]
Transferring through the isomorphism \(C_0[0,\om_1)\cong C([0,\om_1])\),
this proves that the inclusion asked about in \cite{KL2011} is proper for
every nonzero countable ordinal \(\alpha\). This packet does not claim the
\(\alpha=0\) endpoint. The stronger question left open in \cite{KL2014},
whether
\(\overline{\cG}_{c_0(\om_1,C(K_\alpha))}\) equals \(\Sz_{\alpha+1}\), is
also not claimed.

\paragraph{Provenance and novelty check.}
The cheap run indexes had no exact entry for arXiv:1112.4800 or the two
closing ideal-lattice questions. Web searches for exact phrases involving
``Loy--Willis ideal'', ``\(C([0,\omega_1])\)'', ``closed ideals'', and the
properness question mainly rediscovered the source paper. The answer above is
therefore recorded as a literature-implied identification, not as an original
new proof.

\begin{thebibliography}{9}

\bibitem{KL2011}
T. Kania and N. J. Laustsen,
\emph{Uniqueness of the maximal ideal of the Banach algebra of bounded
operators on \(C([0,\omega_1])\)}, arXiv:1112.4800.

\bibitem{KL2014}
T. Kania and N. J. Laustsen,
\emph{Operators on two Banach spaces of continuous functions on locally
compact spaces of ordinals}, arXiv:1304.4951.

\bibitem{KKL}
T. Kania, P. Koszmider, and N. J. Laustsen,
the work cited as \cite{KKL} in \cite{KL2014}, specifically the theorem
characterizing the Loy--Willis ideal as factoring through \(C_0(L_0)\).

\end{thebibliography}

\end{document}
